\documentclass{article}
\usepackage{graphicx}
\usepackage{amsmath}
\usepackage{amssymb}
\usepackage{amsfonts}
\usepackage{hyperref}
\usepackage{exercises}
\usepackage{indentfirst}
\usepackage{circledsteps}
\usepackage{amsthm}
\usepackage{csquotes}
\usepackage{yhmath}
\usepackage{yfonts}
\usepackage{comment}
\usepackage{placeins}
\usepackage{geometry}
\usepackage{cleveref}
\usepackage{stmaryrd}
\usepackage{hyperref}
\usepackage{tikz}
\usepackage{tabto}
\newcommand*{\oldneg}{\mathord{\sim}}
\newcommand{\Z}{\mathbb{Z}}
\newcommand{\N}{\mathbb{N}}
\let\oldemptyset\emptyset
\newcommand{\lbc}{\llbracket}
\newcommand{\rbc}{\rrbracket}
\newcommand{\ltop}[1]{\stackrel{\leftarrow}{#1}}

\theoremstyle{definition}
\newtheorem{definition}{Definition}
\newtheorem{theorem}{Theorem}[section]
\newtheorem{corollary}{Corollary}[theorem]
\newtheorem{lemma}[theorem]{Lemma}

\setlength{\parindent}{.5in}
\geometry{margin=0.5in}
\setlength{\parskip}{12pt}
\title{Chapter 20 Notes}
\begin{document}
\date{}
\maketitle

\section*{Section 2}

\subsection*{Notes}
Note that, in Definition 2.2, $\theta$ is simply being extended; so that $\theta(p)$ is still defined.

It is remarked in Definition 2.1 that, if $\theta$ is an automorphism, then so is $\theta^{-1}$. This is because 
$\theta^{-1}(p)\mathcal{R}\theta^{-1}(q)$ iff $f(g)=p$ and $f(h)=q$ for some $g$ and $q$, meaning 
$f(g)\mathcal{R}f(q)$ iff $g\mathcal{R}q$ (because $f$ is an automorphism).

It is additionaly remarked that, if $\phi$ and $\psi$ are automorphisms, then so is $\phi\psi$. This follows because 
$fg(p)\mathcal{R}fg(q)$ iff $g(p)\mathcal{R}g(q)$ iff $p\mathcal{R}q$.

Note that, if $\theta$ is a function obtained by Definition 2.2, then 
$\theta^{-1}$ is also such a function. That is, 
$\theta^{-1}(f)=\{\langle\theta^{-1}(p),\theta^{-1}(g)\rangle|\langle p,g\rangle\in f\}$ for 
all $f\in\mathcal{D}^\mathcal{G}$ (and the function is still 1-1 and the domain and ranges stay the same). This can be shown in the 
following way. $\theta^{-1}(f)$ is some $h$ such that 
$\{\langle\theta(p),\theta(g)\rangle|\langle p,g\rangle\in h\}=f$. Then 
\begin{center}
\begin{tabular}{ll}
& $\{\langle\theta^{-1}(p),\theta^{-1}(g)\rangle|\langle p,g\rangle\in f\}$ \\
$=$ & $\{\langle\theta^{-1}\theta(p),\theta^{-1}\theta(g)\rangle|\langle p,g\rangle\in h\}$ \\
$=$ & $\{\langle p,g\rangle|\langle p,g\rangle\in h\}$ \\
$=$ & $h$ \\
$=$ & $\theta^{-1}(f)$
\end{tabular}
\end{center}

Also, if $\phi$ and $\psi$ are functions obtained by Definition 2.2, then 
$\phi\psi$ is also such a function. 
That is, $\phi\psi(f)=\{\langle\phi\psi(p),\phi\psi(g)\rangle|\langle p,g\rangle\in f\}$ for all 
$f\in\mathcal{D}^\mathcal{G}$ (this just follows by the definition of $\phi$, $\psi$, and function composition).
The fact that $\phi\psi$ is still 1-1 and has a domain and range of $\mathcal{G}\cup\mathcal{D}^\mathcal{G}$ is also easily seen.

I did not fully understand the paragraph titled \textquote{\textbf{Notation}} on pg. 269. In particular, it was not fully clear to me why 
the recursion taking place in that paragraph was possible.

\subsection*{Exercise 2.1}
We prove Proposition 2.3. First, we show tha
$\theta:\mathcal{D}^\mathcal{G}\xrightarrow[\textrm{onto}]{\text{1-1}}\mathcal{D}^\mathcal{G}$.
By definition, the domain of $\theta$ includes all members of $\mathcal{D}^\mathcal{G}$. We use transfinite induction 
to show that $\theta$ is 1-1, and is such that the range of $\theta$ includes all members and only members of $\mathcal{D}^\mathcal{G}$.

\begin{proof}

The least ordinal case is obvious.

Suppose the result holds for some $\alpha$. That is, for every $f\in\mathcal{R}^\mathcal{G}_\alpha$, the following three things hold:
(1) for all $g\in\mathcal{D}^\mathcal{G}$, if $\theta(f)=\theta(g)$, then $f=g$ (i.e., that $\theta$ is 1-1 for all members 
of $\mathcal{D}^\mathcal{G}$); (2) $f\in\textrm{Ran}(\theta)$; and (3) $\theta(f)\in\mathcal{D}^\mathcal{G}$.
We show that those two conditions also hold for every $f\in\mathcal{R}^\mathcal{G}_{\alpha+1}$. As for (1), suppose that 
$f\in\mathcal{R}^\mathcal{G}_{\alpha+1}$ and, for some $a$, $\theta(f)=\theta(a)$. 
Now, $\theta(f)=\{\langle\theta(p),\theta(g)\rangle|\langle p,g\rangle\in f\}$ and 
$\theta(a)=\{\langle\theta(p),\theta(g)\rangle|\langle p,g\rangle\in a\}$.
To show that $f=a$, we show that $f\subseteq a$ and $a\subseteq f$. We prove the first case, that is, $f\subseteq a$; the other case 
is similar. Suppose $x\in f$. Then $x=\langle p,a\rangle$ for some $p$ and $a$. Then it follows that 
$\langle\theta(p),\theta(g)\rangle\in\theta(f)$. But, because we assumed $\theta(f)=\theta(a)$, it follows that 
$\langle\theta(p),\theta(g)\rangle\in\theta(a)$. Hence, by the definition of $\theta(a)$, it follows that 
$\langle q,h\rangle\in a$ for some $q,h$ where $\theta(p)=\theta(q)$ and $\theta(g)=\theta(h)$. But, by the induction assumption, the 
fact that $g\in\mathcal{R}^\mathcal{G}_\alpha$ and the fact that $\theta\restriction\mathcal{G}$ is 1-1, it follows that 
$p=q$ and $g=h$. Thus $\langle p,g\rangle=\langle q,h\rangle$, meaning $x\in a$.

Next we show that (2): for every $f\in\mathcal{R}^\mathcal{G}_{\alpha+1}$, $f\in\textrm{Ran}(\theta)$. Now, all members of 
$f$ are of the form $\langle p,g\rangle$ where $g\in\mathcal{R}^\mathcal{G}_\alpha$. Hence, it follows by the induction assumption 
that $g\in\textrm{Ran}(\theta)$ for each such $g$. Define a set $k$ as follows:
$k=\{\langle q,h\rangle|\theta(q)=p\textrm{ and }\theta(h)=g\textrm{ for some }\langle p,g\rangle\in f\}$. 
Such a set exists by the induction assumption, and because $k$ is a set, it follows that it must be a subset of 
$R^\mathcal{G}_\beta$ for some $\beta$, and consequently, it is a member of $R^\mathcal{G}_{\beta+1}$ (I was not fully 
sure of this, though it seemed to be true). Hence $k$ is in the domain of $\theta$, and because 
$\theta(k)=f$, it follows that $f\in\textrm{Ran}(\theta)$.

Lastly, we show that (3): for every $f\in\mathcal{R}^\mathcal{G}_{\alpha+1}$, $\theta(f)\in\mathcal{D}^\mathcal{G}$.
By definition, $\theta(f)=\{\langle\theta(p),\theta(g)\rangle|\langle p,g\rangle\in f\}$. By the induction assumption, and 
the fact that $\theta\restriction\mathcal{G}$ is 1-1 onto $\mathcal{G}$, it follows that 
$\theta(f)\subseteq R^\mathcal{G}_\beta$ for some $\beta$, and consequently, it is a member of $R^\mathcal{G}_{\beta+1}$ (again, I was 
not fully sure of this). Hence $\theta(f)\in\mathcal{D}^\mathcal{G}$. 

The limit ordinal case is obvious.

\end{proof}

Thus we know, by transfinite induction, that $\theta$ is 1-1, and is such that the range of $\theta$ includes all members and only members of 
$\mathcal{D}^\mathcal{G}$. Hence we know $\theta:\mathcal{D}^\mathcal{G}\xrightarrow[\textrm{onto}]{\text{1-1}}\mathcal{D}^\mathcal{G}$. The last 
two things we must show are that, if $e$ is the identity automorphism of $\langle\mathcal{G},\mathcal{R}\rangle$, then $e$ is the identity 
map on $\mathcal{D}^\mathcal{G}$, and similarly for inverses.

First, suppose that $e$ is the identity automorphism of $\langle\mathcal{G},\mathcal{R}\rangle$. We prove, by transfinite induction, 
that $e(f)=f$ for all $f\in\mathcal{D}^\mathcal{G}$.

\begin{proof}

The least ordinal case is obvious, and likewise for the limit ordinal case.

Suppose the result holds for some $\alpha$, that is, for all $f\in R^\mathcal{G}_\alpha$, $e(f)=f$. Consider any 
$f\in R^\mathcal{G}_{\alpha+1}$. Then, if $f\ne\emptyset$, 
\begin{center}
\begin{tabular}{llll}
$e(f)$ & $=$ & $\{\langle e(p),e(g)\rangle|\langle p,g\rangle\in f\}$ \\
& $=$ & $\{\langle p,g\rangle|\langle p,g\rangle\in f\}$ & (By induction assumption and fact that $e$ is identity map) \\
& $=$ & $f$
\end{tabular}
\end{center}
On the other hand, if $f=\emptyset$, $e(f)=f$ is obvious, and this completes the proof.

\end{proof}

Second, suppose that $\theta^{-1}$ is an inverse of some automorphism $\theta$ of 
$\langle\mathcal{G},\mathcal{R}\rangle$. We show that $\theta^{-1}$ is the inverse of 
$\theta$ on $\mathcal{D}^\mathcal{G}$. We prove this by showing that, for all $f\in\mathcal{D}^\mathcal{G}$, 
$\theta\theta^{-1}(f)=f$ and $\theta^{-1}\theta=f$. We do this by using transfinite induction.

\begin{proof}

Again, the least ordinal and limit ordinal cases are obvious.

Suppose the result holds for some $\alpha$, that is, for all $f\in R^\mathcal{G}_\alpha$, 
$\theta\theta^{-1}(f)=f$ and $\theta^{-1}\theta=f$. Consider any 
$f\in R^\mathcal{G}_{\alpha+1}$ (where $f\ne\emptyset$; the case where $f=\emptyset$ is obvious). 
We first show that $\theta\theta^{-1}f=f$. Now, 
$\theta^{-1}(f)=\{\langle\theta^{-1}(p),\theta^{-1}(g)\rangle|\langle p,g\rangle\in f\}$. Hence
\begin{center}
\begin{tabular}{ll}
& $\theta\theta^{-1}(f)$ \\
$=$ & $\{\langle\theta\theta^{-1}(p),\theta\theta^{-1}(g)\rangle|\langle p,g\rangle\in f\}$ \\
$=$ & $\{\langle p,q\rangle|\langle p,g\rangle\in f\}$ \\
$=$ & $f$
\end{tabular}
\end{center}
by the induction assumption and the fact that $\theta^{-1}$ is the inverse of $\theta$ on $\langle\mathcal{G},\mathcal{R}\rangle$.
Thus we know that $\theta\theta^{-1}(f)=f$. Next we must show that $\theta^{-1}\theta=f$.
Well, $\theta(f)=\{\langle\theta(p),\theta(g)\rangle|\langle p,g\rangle\in f\}$, and thus
\begin{center}
\begin{tabular}{ll}
& $\theta^{-1}\theta(f)$ \\
$=$ & $\{\langle\theta^{-1}\theta(p),\theta^{-1}\theta(g)\rangle|\langle p,g\rangle\in f\}$ \\
$=$ & $\{\langle p,q\rangle|\langle p,g\rangle\in f\}$ \\
$=$ & $f$
\end{tabular}
\end{center} 
by the induction assumption.

\end{proof}

This completes the proof of Proposition 2.3.

\subsection*{Exercise 2.2}
We prove that if $a$ is an $M$-set (i.e., $a\in M$), $\theta(\hat{a})=\hat{a}$ for every automorphism $\theta$. Then we show that 
the counterparts of $M$-sets (i.e., $\hat{a}$ for each $a\in M$) are \textfrak{J}-invariant for every automorphism sequence \textfrak{G}. 
First we prove by transfinite induction (on the $\alpha$ for which $\hat{a}\in R^\mathcal{G}_\alpha$)
that $\theta(\hat{a})=\hat{a}$ for every automorphism $\theta$ and every $M$-set $a$.

\begin{proof}
Consider any automorphism $\theta$. The least ordinal case is vacuous, and the limit ordinal case is obvious.

Suppose the result holds for some $\alpha$. That is, suppose that for every $\hat{a}\in R^\mathcal{G}_\alpha$, 
$\theta(\hat{a})=\hat{a}$. Consider any $\hat{b}\in R^\mathcal{G}_{\alpha+1}$. If $\hat{b}=\emptyset$, then 
the result obviously follows. So suppose $\hat{b}\ne\emptyset$. Then 
$\hat{b}=\{\langle p,\hat{g}\rangle|p\in\mathcal{G}\textrm{ and }g\in b\}$. Hence 
\begin{center}
\begin{tabular}{lll}
& $\theta(\hat{b})=\{\langle\theta(p),\theta(\hat{g})\rangle|\langle p,g\rangle\in b\}$ \\
$=$ & $\{\langle\theta(p),\hat{g}\rangle|\langle p,g\rangle\in b\}$ & (By the induction assumption) \\
$=$ & $\hat{b}$ & (Because $\theta:\mathcal{G}\xrightarrow[\textrm{onto}]{\text{1-1}}\mathcal{G}$)
\end{tabular}
\end{center}

\end{proof}

Now consider any automorphism sequence \textfrak{J} and some $M$-set $a$. Then, by the above result, 
$\theta(\hat{a})=\hat{a}$ for every automorphism $\theta$. Consequently 
$\hat{a}$ is \textfrak{G}-invariant for every automorphism group \textfrak{G} of $\langle\mathcal{G},\mathcal{R}\rangle$. To show that 
$\hat{a}$ is \textfrak{J}-invariant, we must show that (1) for some \textfrak{G} in the sequence 
\textfrak{J}, $\hat{a}$ is \textfrak{G}-invariant; (2) for each $\langle p,g\rangle\in\hat{a}$, $g$ is 
\textfrak{J}-invariant. We prove these two things by induction on the $\mathcal{G}$-rank of $\hat{a}$.

\begin{proof}
Again, the least ordinal case is vacuous. Also, the case when 
$\hat{a}=\emptyset$ is likewise obvious, so suppose $\hat{a}\ne\emptyset$. Now suppose the result holds for all ordinals less than some 
ordinal $\alpha$. We show that, for all $\hat{a}$ with $\mathcal{G}$-rank $\alpha$, both (1) and (2) hold.

(1) This holds because, as remarked just before the start of this proof, 
$\hat{a}$ is \textfrak{G}-invariant for \textit{every} automorphism group \textfrak{G} of $\langle\mathcal{G},\mathcal{R}\rangle$,
hence also \textfrak{G}-invariant for every \textfrak{G} in the sequence \textfrak{J}.

(2) Consider any $\langle p,g\rangle\in\hat{a}$. By definition of $\hat{a}$, it follows that 
$g=\hat{b}$ for some $b$, and $\hat{b}$ must have $\mathcal{G}$-rank less than $\alpha$. Thus, by the induction 
assumption, $\hat{b}$, that is, $g$, must be \textfrak{J}-invariant.

\end{proof}

This completes the proof of Proposition 2.9.

\section*{Section 3}

\subsection*{Theorem 3.3}
We show that Theorem 3.3 is provable from Lemma 3.4.

\begin{theorem}
Suppose \textfrak{J} is an automorphism sequence of $\langle\mathcal{G},\mathcal{R}\rangle$ and $\theta\in\textfrak{G}_0$. 
If $f$ is \textfrak{J}-invariant, so is $\theta(f)$.
\end{theorem}

\begin{proof}
By transfinite induction on the $\alpha$ for which $f\in R^\mathcal{G}_\alpha$. 
The least ordinal and limit ordinal cases are obvious. Suppose the result holds for some $\alpha$. If 
$f=\emptyset$, then the result easily follows (because $\theta(f)=f$). So suppose $f\ne\emptyset$. Because 
$f$ is \textfrak{J}-invariant, it follows by definition that (1) for some 
$\textfrak{G}_k$ in the sequence \textfrak{J}, $f$ is $\textfrak{G}_k$-invariant; (2) for each 
$\langle p,g\rangle\in f$, $g$ is \textfrak{J}-invariant. To prove that 
$\theta(f)$ is \textfrak{J}-invariant, we must show that (3) for some 
$\textfrak{G}_j$ in the sequence \textfrak{J}, $f$ is $\textfrak{G}_j$-invariant; (2) for each
$\langle\theta(p),\theta(g)\rangle\in\theta(f)$ (for $\langle p,g\rangle\in f$), $\theta(g)$ is 
\textfrak{J}-invariant. By Lemma 3.4, and the fact that $f$ is $\textfrak{G}_k$-invariant, 
$\theta(f)$ is also $\textfrak{G}_k$-invariant, from which (3) follows immediately. As for 
(4), consider any $\langle\theta(p),\theta(g)\rangle\in\theta(f)$ where $\langle p,g\rangle\in f$. Then, by 
(2), $g$ is \textfrak{J}-invariant. But $g\in R^\mathcal{G}_\alpha$, and thus, by the induction assumption, 
$\theta(g)$ is also \textfrak{J}-invariant. This completes the proof.

\end{proof}

\subsection*{Theorem 3.5}
We provide details of the proof of Theorem 3.5.

\begin{theorem}
Let $X$ be any closed formula of $\mathcal{L}^*_C$ (the language with only $\in$ as a relation symbol) having 
all constants in $\mathcal{D}^\mathcal{G}_\textfrak{J}$. For any $p\in\mathcal{G}$ and any $\theta\in\textfrak{G}_0$, 
$p\vDash_\textfrak{J} X$ iff $\theta(p)\vDash_\textfrak{J}\theta(X)$.
\end{theorem}

\begin{proof}
First we note the things remarked in proof given of the theorem in the book. By the previous theorem, if $X$ has all its 
constants in $\mathcal{D}^\mathcal{G}_\textfrak{J}$, the same is true for $\theta(X)$, so 
$\theta(p)\vDash_\textfrak{J}\theta(X)$ is a sensical statement. Now, the result is true for atomic $X$ in 
$\mathcal{L}^*_C$, since at the atomic level, $\vDash_\textfrak{J}$ and $\vDash$ agree, and we 
have Theorem 3.2. Now we prove the result for non-atomic formulas by induction on degree. This is done by considering 
the negation, disjunction, existential quantifier and the modal necessity cases.

Negation case: Suppose the result holds for $X$. We show it also holds for $\oldneg X$. That is, we show that 
$p\vDash_\textfrak{J}\oldneg X$ iff $\theta(p)\vDash_\textfrak{J}\theta(\oldneg X)$. On the one hand, suppose that 
$p\vDash_\textfrak{J}\oldneg X$. Then $p\nvDash_\textfrak{J} X$, meaning, by the induction assumption, 
$\theta(p)\nvDash_\textfrak{J}\theta(X)$, hence 
$\theta(p)\vDash_\textfrak{J}\oldneg\theta(X)$, which is the same as saying 
$\theta(p)\vDash_\textfrak{J}\theta(\oldneg X)$. On the other hand, suppose that 
$\theta(p)\vDash_\textfrak{J}\theta(\oldneg X)$. This is the same as saying 
$\theta(p)\vDash_\textfrak{J}\oldneg\theta(X)$, hence 
$\theta(p)\nvDash_\textfrak{J}\theta(X)$, meaning, by the induction assumption,
$p\nvDash X$, hence $p\vDash\oldneg X$. Thus the negation case is complete.

Disjunction case: Suppose the result holds for $X$ and $Y$. We show it also holds for $X\vee Y$. That is, we show that 
$p\vDash_\textfrak{J} X\vee Y$ iff $\theta(p)\vDash_\textfrak{J}\theta(X\vee Y)$. On the one hand, suppose that 
$p\vDash_\textfrak{J} X\vee Y$. If $\theta(p)\nvDash_\textfrak{J}\theta(X\vee Y)$, we would have 
$\theta(p)\nvDash_\textfrak{J}\theta(X)\vee\theta(Y)$, hence both 
$\theta(p)\nvDash_\textfrak{J}\theta(X)$ and $\theta(p)\nvDash_\textfrak{J}\theta(Y)$, meaning, by the induction assumption, 
both $p\nvDash_\textfrak{J} X$ and $p\nvDash_\textfrak{J} Y$, meaning $p\vDash_\textfrak{J}\oldneg (X\vee Y)$, a contradiction. Hence we must have 
$\theta(p)\vDash_\textfrak{J}\theta(X\vee Y)$. On the other hand, suppose
$\theta(p)\vDash_\textfrak{J}\theta(X\vee Y)$. Then, if $p\nvDash X\vee Y$, we would have 
both $p\nvDash X$ and $p\nvDash Y$, meaning, by the induction assumption, both 
$\theta(p)\nvDash_\textfrak{J}\theta(X)$ and $\theta(p)\nvDash_\textfrak{J}\theta(Y)$, hence
$\theta(p)\vDash_\textfrak{J}\oldneg(\theta(X)\vee\theta(Y))$, that is, 
$\theta(p)\vDash_\textfrak{J}\oldneg\theta(X\vee Y)$, a contradiction. Thus it must be that 
$p\vDash X\vee Y$. This concludes the disjunction case.

Existential quantifier case: Suppose, for all $a\in\mathcal{D}^\mathcal{G}_\textfrak{J}$, 
the result holds for $\phi(a)$ (for some $\phi(x)$ with $x$ being the only free variable of $\phi$). We must show that 
result also holds for $(\exists x)\phi(x)$. That is, we show that 
$p\vDash_\textfrak{J}(\exists x)\phi(x)$ iff $\theta(p)\vDash_\textfrak{J}\theta((\exists x)\phi(x))$. On the one hand, suppose 
$p\vDash_\textfrak{J}(\exists x)\phi(x)$. Then, for some $a\in\mathcal{D}^\mathcal{G}_\textfrak{J}$, 
$p\vDash_\textfrak{J}\phi(a)$. Then, by the induction assumption, 
$\theta(p)\vDash_\textfrak{J}\theta(\phi(a))$. Let $\theta((\exists x)\phi(x))=(\exists x)\psi(x)$. 
Now, $\theta(\phi(a))=\psi(\theta(a))$, meaning $\theta(p)\vDash_\textfrak{J}\psi(\theta(a))$, and thus 
$\theta(p)\vDash_\textfrak{J}(\exists x)\psi(x)$, that is, $\theta(p)\vDash_\textfrak{J}\theta((\exists x)\phi(x))$.
On the other hand, suppose $\theta(p)\vDash_\textfrak{J}\theta((\exists x)\phi(x))$. Again, let 
$\theta((\exists x)\phi(x))=(\exists x)\psi(x)$. Thus we have 
$\theta(p)\vDash_\textfrak{J}(\exists x)\psi(x)$, meaning $\theta(p)\vDash_\textfrak{J}\psi(a)$ for some $a$. Now, 
$\psi(a)=\theta(\phi(\theta^{-1}(a)))$, meaning $\theta(p)\vDash_\textfrak{J}\theta(\phi(\theta^{-1}(a)))$, and thus, 
by the induction assumption, $p\vDash_\textfrak{J}\phi(\theta^{-1}(a))$, meaning 
$p\vDash_\textfrak{J}(\exists x)\phi(x)$. This completes the existential quantifier case.

Modal necessity case: Suppose the result holds for $X$. We must show it also holds for $\Box X$. That is, we show that 
$p\vDash_\textfrak{J}\Box X$ iff $\theta(p)\vDash_\textfrak{J}\theta(\Box X)$, which is just saying
$\theta(p)\vDash_\textfrak{J}\Box\theta(X)$. On the one hand, suppose $p\vDash_\textfrak{J}\Box X$. To show 
that $\theta(p)\vDash_\textfrak{J}\Box\theta(X)$, consider any $q$ with 
$\theta(p)\mathcal{R}q$. Then, because $\theta:\mathcal{G}\xrightarrow[\text{onto}]{\text{1-1}}\mathcal{G}$, 
$\theta(p')=q$ for some $p'$. Hence $\theta(p)\mathcal{R}\theta(p')$, meaning, because $\theta$ is an 
automorphism, $p\mathcal{R}p'$, and thus $p'\vDash_\textfrak{J} X$, which implies 
$\theta(p')\vDash_\textfrak{J}\theta(X)$ by the induction assumption, that is, 
$q\vDash_\textfrak{J}\theta(X)$. Because $q$ was arbitrary, it follows that 
$\theta(p)\vDash_\textfrak{J}\Box\theta(X)$. 
On the other hand, suppose $\theta(p)\vDash_\textfrak{J}\Box\theta(X)$. To show that 
$p\vDash_\textfrak{J}\Box X$, consider any $q$ with $p\mathcal{R}q$. Because $\theta$ is an 
automorphism, it follows that $\theta(p)\mathcal{R}\theta(q)$, and thus
$\theta(q)\vDash_\textfrak{J}\theta(X)$, meaning, by the induction assumption, 
$q\vDash_\textfrak{J} X$. Because $q$ was arbitrary, it follows that 
$p\vDash_\textfrak{J}\Box X$. This concludes the modal necessity case, therefore completing the proof.

\end{proof}

\section*{Section 4}

\subsection*{Proposition 4.1}
It is remakred in the second sentence of the first paragraph of the proof of Proposition 4.1 
that $f\in g$ and $\lbc f\in g\rbc$ are equivalent. This is 
not correct; $p\vDash f\in g$ implies $p\vDash\lbc f\in g\rbc$, but $p\vDash\lbc f\in g\rbc$ does not necessarily imply
$p\vDash f\in g$. However, the atomic case is still simple as $p\vDash f\in g$ iff 
$p\vDash_\textfrak{J} f\in g$ for every $p$ (and $f,g\in\mathcal{D}^\mathcal{G}_\textfrak{J}$), by the definition of
$\vDash_\textfrak{J}$. It follows that $p\vDash\Box\Diamond(f\in g)$ iff 
$p\vDash_\textfrak{J}\Box\Diamond(f\in g)$ for every $p$, that is, 
$p\vDash\lbc f\in g\rbc$ iff $p\vDash_\textfrak{J}\lbc f\in g\rbc$ for every $p$.

As a second note, we provide a proof of the second case mentioned, but omitted, in the proof given in the book. 

\begin{proof}
That is, we suppose 
$\phi(x_1,...,x_n)$ is $(\exists y)[y\in x_1\wedge\psi(y,x_1,...,x_n)]$, the result is known for $\psi$, 
$p\vDash_\textfrak{J}\lbc\phi(s_1,...,s_n)\rbc$, where $s_1,...,s_n\in\mathcal{D}^\mathcal{G}_\textfrak{J}$, but 
$p\nvDash\lbc\phi(s_1,...,s_n)\rbc$. We derive a contradiction.

Because $p\nvDash\lbc\phi(s_1,...,s_n)\rbc$, $p'\vDash\lbc\oldneg\phi(s_1,...,s_n)\rbc$ for some 
$p'$ with $p\mathcal{R}p'$, or equivalently,
$p'\vDash\lbc(\forall y)[y\in s_1\supset\oldneg\psi(y,s_1,...,s_n)]\rbc$.

Since $p\vDash_\textfrak{J}\lbc\phi(s_1,...,s_n)\rbc$, also 
$p'\vDash_\textfrak{J}\lbc\phi(s_1,...,s_n)\rbc$, that is, 
$p'\vDash_\textfrak{J}\lbc(\exists y)[y\in s_1\wedge\psi(y,s_1,...,s_n)]\rbc$. Then for some 
$p''$ with $p'\mathcal{R}p''$ and some $a\in\mathcal{D}^\mathcal{G}_\textfrak{J}$,
$p''\vDash_\textfrak{J}\lbc[a\in s_1\wedge\psi(a,s_1,...,s_n)]\rbc$. It follows that both 
$p''\vDash_\textfrak{J}\lbc (a\in s_1)\rbc$ and $p''\vDash_\textfrak{J}\lbc\psi(a,s_1,...,s_n)\rbc$. Because the atomic case has been established, 
$p''\vDash\lbc (a\in s_1)\rbc$. But also, since $p'\mathcal{R}p''$,
$p''\vDash\lbc(\forall y)[y\in s_1\supset\oldneg\psi(y,s_1,...,s_n)]\rbc$, which implies
$p''\vDash\lbc\oldneg\psi(a,s_1,...,s_n)\rbc$. But 
$p''\vDash_\textfrak{J}\lbc\psi(a,s_1,...,s_n)\rbc$ has also been established, and this contradicts the induction
hypothesis.

\end{proof}

\subsection*{Exercise 4.1}

(a) We suppose that $X$ is $\Pi$-formula, and show that $\oldneg X$ is equivalent to some 
$\Sigma$-formula.

\begin{proof}
By induction on the complexity of $X$. The atomic case is obvious (because atomic formulas are $\Delta_0$, hence both 
$\Sigma$ and $\Pi$ formulas). The different cases are for disjunction, conjunction, universal quantification, 
bounded universal quantification and bounded existential quantification.

Disjunction case: Suppose the result holds for two $\Pi$-formulas $X$ and $Y$. We must show that it also holds for 
$X\vee Y$. Well, $\oldneg(X\vee Y)$ is equivalent to $\oldneg X\wedge\oldneg Y$. By the induction assumption, 
$\oldneg X$ and $\oldneg Y$ are equivalent to two $\Sigma$-formulas $X'$ and $Y'$, respectively. Hence 
$\oldneg(X\vee Y)$ is equivalent to $X'\wedge Y'$, which is a $\Sigma$-formula. This completes the disjunction case.

Conjunction case: Suppose the result holds for two $\Pi$-formulas $X$ and $Y$. We must show that it also holds for 
$X\wedge Y$. Well, $\oldneg(X\wedge Y)$ is equivalent to $\oldneg X\vee\oldneg Y$. By the induction assumption,
$\oldneg X$ and $\oldneg Y$ are equivalent to two $\Sigma$-formulas $X'$ and $Y'$, respectively. Hence
$\oldneg(X\wedge Y)$ is equivalent to $X'\vee Y'$, which is a $\Sigma$-formula. This completes the conjunction case.

Universal quantification case: Suppose the result holds for some $\Pi$-formula $\phi$. We must show it also holds for 
$(\forall x)\phi$. Well, $\oldneg(\forall x)\phi$ is equivalent to 
$(\exists x)\oldneg\phi$. By the induction assumption, $\oldneg\phi$ is equivalent to $\psi$ for some 
$\Sigma$-formula $\psi$. Hence $\oldneg(\forall x)\phi$ is equivalent to 
$(\exists x)\psi$, which is a $\Sigma$-formula. This completes the universal quantification case.

Bounded universal quantification: Suppose the result holds for some $\Pi$-formula $\phi$. We must show it 
also holds for $(\forall x\in y)\phi$. Well, $\oldneg(\forall x\in y)\phi$ is 
equivalent to $(\exists x\in y)\oldneg\phi$. By the induction assumption, $\oldneg\phi$ is equivalent to 
$\psi$ for some $\Sigma$-formula $\psi$. Hence $\oldneg(\forall x\in y)\phi$ is equivalent to
$(\exists x\in y)\psi$, which is a $\Sigma$-formula. This completes the bounded universal quantification case.

Bounded existential quantification: Suppose the result holds for some $\Pi$-formula $\phi$. We must show it 
also holds for $(\exists x\in y)\phi$. Well, $\oldneg(\exists x\in y)\phi$ is 
equivalent to $(\forall x\in y)\oldneg\phi$. By the induction assumption, $\oldneg\phi$ is equivalent to
$\psi$ for some $\Sigma$-formula $\psi$. Hence $\oldneg(\exists x\in y)\phi$ is equivalent to
$(\forall x\in y)\psi$, which is a $\Sigma$-formula. This completes the bounded existential quantification case, 
therefore completing the proof of (a).

\end{proof}

(b) The proof of (b) is nearly identical to that of (a), so we omit it.

\subsection*{Exercise 4.2}
We prove Corollary 4.3. That is, let $\phi(x_1,...,x_n)$ be a formula of $\mathcal{L}_C$, and  $s_1,...,s_n$ be members of 
$\mathcal{D}^\mathcal{G}_\textfrak{J}$. If $\phi$ is a $\Sigma$-formula and 
$p\vDash_\textfrak{J}\lbc\phi(s_1,...,s_n)\rbc$, we show that $p\vDash\lbc\phi(s_1,...,s_n)\rbc$. If $\phi$ is a 
$\Pi$-formula and $p\vDash\lbc\phi(s_1,...,s_n)\rbc$, we show that $p\vDash_\textfrak{J}\lbc\phi(s_1,...,s_n)\rbc$.

\begin{proof}
By induction on the complexity of $\phi$. 

(a) First, suppose that $\phi$ is a $\Sigma$-formula. Also suppose that 
$p\vDash_\textfrak{J}\lbc\phi(s_1,...,s_n)\rbc$. We must show that $p\vDash\lbc(s_1,...,s_n)\rbc$.
The case where $\phi$ is also a $\Delta_0$ formula is immediate 
from Proposition 4.1. Thus we need only consider the disjunction, conjunction, existential quantifier, 
bounded universal quantifier and bounded existential quantifier cases. 

Disjunction case: Suppose $\phi(s_1,...,s_n)$ is of the form $\psi(s_1,...,s_n)\vee\delta(s_1,...,s_n)$, but
$p\nvDash\lbc\phi(s_1,...,s_n)\rbc$. We derive a contradiction. Because 
$p\nvDash\lbc\phi(s_1,...,s_n)\rbc$, it follows that 
$p'\vDash\lbc\oldneg\phi(s_1,...,s_n)\rbc$ for some $p'$ with $p\mathcal{R}p'$. Also, for 
some $p''$ with $p'\mathcal{R}p''$, at least one of $p''\vDash_\textfrak{J}\lbc\psi(s_1,...,s_n)\rbc$ and 
$p''\vDash_\textfrak{J}\lbc\delta(s_1,...,s_n)\rbc$ holds. Suppose, without loss of generality, that the former is the case, that is, 
$p''\vDash_\textfrak{J}\lbc\psi(s_1,...,s_n)\rbc$. By the induction assumption, this means 
$p''\vDash\lbc\psi(s_1,...,s_n)\rbc$. Because $p'\mathcal{R}p''$, $p''\vDash\lbc\oldneg\phi(s_1,...,s_n)\rbc$, 
and hence $p''\nvDash\lbc\phi(s_1,...,s_n)\rbc$, meaning $p''\nvDash\lbc\psi(s_1,...,s_n)\rbc$, a contradiction.

Conjunction case: Suppose $\phi(s_1,...,s_n)$ is of the form $\psi(s_1,...,s_n)\wedge\delta(s_1,...,s_n)$, but
$p\nvDash\lbc\phi(s_1,...,s_n)\rbc$. We derive a contradiction. Because
$p\nvDash\lbc\phi(s_1,...,s_n)\rbc$, it follows that 
$p'\vDash\lbc\oldneg\phi(s_1,...,s_n)\rbc$ for some $p'$ with $p\mathcal{R}p'$. Also,
$p'\vDash_\textfrak{J}\lbc\psi(s_1,...,s_n)\rbc$ and 
$p'\vDash_\textfrak{J}\lbc\delta(s_1,...,s_n)\rbc$ both hold, meaning, by the induction assumption, 
both $p'\vDash\lbc\psi(s_1,...,s_n)\rbc$ and 
$p'\vDash\lbc\delta(s_1,...,s_n)\rbc$ hold. Because 
$p'\vDash\lbc\oldneg\phi(s_1,...,s_n)\rbc$, either 
$p'\nvDash\lbc\psi(s_1,...,s_n)\rbc$ or $p'\nvDash\lbc\delta(s_1,...,s_n)\rbc$, and this is a contradiction.

Existential quantifier case: Suppose $\phi(s_1,...,s_n)$ is of the form 
$(\exists x)\phi(x,s_1,...,s_n)$, but $p\nvDash\lbc(\exists x)\phi(x,s_1,...,s_n)\rbc$. We derive a contradiction. 
This means $p'\vDash\lbc\oldneg(\exists x)\phi(x,s_1,...,s_n)\rbc$ for some $p'$ with $p\mathcal{R}p'$. Because 
$p\vDash_\textfrak{J}\lbc(\exists x)\phi(x,s_1,...,s_n)\rbc$, it follows by the induction assumption that, for some 
$f\in\mathcal{D}^\mathcal{G}_\textfrak{J}$ and some $p''$ with $p'\mathcal{R}p''$,
$p''\vDash\lbc\phi(f,s_1,...,s_n)\rbc$. Then $p''\vDash\lbc(\exists x)(x,s_1,...s_n)\rbc$. But 
we also have $p'\vDash\lbc\oldneg(\exists x)\phi(x,s_1,...,s_n)\rbc$, and this implies
$p''\vDash\lbc\oldneg(\exists x)\phi(x,s_1,...,s_n)\rbc$, a contradiction.

Bounded universal quantifier case: Suppose $\phi(s_1,...,s_n)$ is of the form
$(\forall x\in s_1)\psi(x,s_1,...,s_n)$, but $p\nvDash\lbc\phi(s_1,...,s_n)\rbc$. We derive a contradiction.
This means $p'\vDash\lbc\oldneg(\forall x\in s_1)\phi(x,s_1,...,s_n)\rbc$ for some $p'$ with $p\mathcal{R}p'$, which 
implies $p''\vDash\lbc d\in s_1\rbc$ and $p''\vDash\lbc\oldneg\phi(d,s_1,...,s_n)\rbc$ for some 
$p''$ with $p'\mathcal{R}p''$ and some $d\in\mathcal{D}^\mathcal{G}$. Because 
$p''\vDash\lbc d\in s_1\rbc$, it follows that $p'''\vDash\lbc f\epsilon s_1\rbc$ and 
$p'''\vDash\lbc d\approx f\rbc$ for some $f$ and some $p'''$ with $p''\mathcal{R}p'''$. 
Because $s_1\in\mathcal{D}^\mathcal{G}_\textfrak{J}$, 
$s_1$ is \textfrak{J}-invariant, which means, by Proposition 2.8, that $g$ is \textfrak{J}-invariant as well, and thus 
$g\in\mathcal{D}^\mathcal{G}_\textfrak{J}$. Then, by substitutivity, 
$p'''\vDash\lbc\oldneg\phi(g,s_1,...,s_n)\rbc$. Also, by Proposition 4.1, 
$p'''\vDash_\textfrak{J}\lbc g\in s_1\rbc$. But also 
$p'''\vDash_\textfrak{J}\lbc(\forall x\in s_1)\psi(x,s_1,...,s_n)\rbc$ (because 
$p\vDash_\textfrak{J}\lbc(\forall x\in s_1)\psi(x,s_1,...,s_n)\rbc$), which means
$p'''\vDash_\textfrak{J}\lbc\psi(g,s_1,...,s_n)\rbc$, hence
$p'''\vDash\lbc\psi(g,s_1,...,s_n)\rbc$ by the induction hypothesis, which is a contradiction.

Bounded existential quantifier case: This is even easier than the bounded universal quantifier case, and we omit it.

(b) Next, suppose that $\phi$ is a $\Pi$-formula. Also suppose that $p\vDash\lbc\phi(s_1,...,s_n)\rbc$. We 
must show that $p\vDash_\textfrak{J}\lbc\phi(s_1,...,s_n)\rbc$. The case where $\phi$ is also a $\Delta_0$ formula
is immediate from Proposition 4.1. The disjunction and conjunction cases are identical to those of part (a), 
while the bounded universal quantifier case is even easier than that of (a). Thus we only consider the 
universal quantifier and bounded existential quantifier cases.

Universal quantifier case: Suppose $\phi(s_1,...,s_n)$ is of the form
$(\forall x)\phi(x,s_1,...,s_n)$, but $p\nvDash_\textfrak{J}\lbc(\forall x)\phi(x,s_1,...,s_n)\rbc$. We derive a contradiction.
This means $p'\vDash_\textfrak{J}\lbc\oldneg\phi(d,s_1,...,s_n)\rbc$ for some $d\in\mathcal{D}^\mathcal{G}_\textfrak{J}$ and 
some $p'$ with $p\mathcal{R}p'$. Because $p\vDash\lbc(\forall x)\phi(x,s_1,...,s_n)\rbc$,
also $p'\vDash\lbc(\forall x)\phi(x,s_1,...,s_n)\rbc$, meaning
$p'\vDash\lbc\phi(d,s_1,...,s_n)\rbc$. But by the induction hypothesis, this means
$p'\vDash_\textfrak{J}\lbc\phi(d,s_1,...,s_n)\rbc$, and this is a contradiction.

Bounded existential quantifier case: Suppose $\phi(s_1,...,s_n)$ is of the form
$(\exists x\in s_1)\phi(x,s_1,...,s_n)$, but $p\nvDash_\textfrak{J}\lbc(\exists x\in s_1)\phi(x,s_1,...,s_n)\rbc$. We derive a contradiction.
Because $p\nvDash_\textfrak{J}\lbc(\exists x\in s_1)\phi(x,s_1,...,s_n)\rbc$, it follows that
$p'\vDash_\textfrak{J}\lbc\oldneg(\exists x\in s_1)\phi(x,s_1,...,s_n)\rbc$ for some $p'$ with $p\mathcal{R}p'$, that is, 
$p'\vDash_\textfrak{J}\lbc(\forall x\in s_1)\oldneg\phi(x,s_1,...,s_n)\rbc$. Because 
$p\vDash\lbc(\exists x\in s_1)\phi(x,s_1,...,s_n)\rbc$, also
$p'\vDash\lbc(\exists x\in s_1)\phi(x,s_1,...,s_n)\rbc$, and thus
$p''\vDash\lbc d\in s_1\rbc$ and $p''\vDash\lbc\phi(d,s_1,...,s_n)\rbc$ for some 
$d\in\mathcal{D}^\mathcal{G}$. Because $p''\vDash\lbc d\in s_1\rbc$, 
$p'''\vDash\lbc d\approx f\rbc$ and $p'''\vDash\lbc f\epsilon s_1\rbc$ both hold for some $f$. Using similar reasoning as in the 
bounded universal quantifier case of (a), it follows that $f\in\mathcal{D}^\mathcal{G}_\textfrak{J}$. By substitutivity,
$p'''\vDash\lbc\phi(f,s_1,...,s_n)\rbc$, meaning
$p'''\vDash_\textfrak{J}\lbc\phi(f,s_1,...,s_n)\rbc$ by the induction hypothesis. But also
$p'''\vDash_\textfrak{J}\lbc f\in s_1\rbc$ (by Proposition 4.1) and 
$p'''\vDash_\textfrak{J}\lbc(\forall x\in s_1)\oldneg\phi(x,s_1,...,s_n)\rbc$ (because
$p'\mathcal{R}p'''$), meaning $p'''\vDash\lbc\oldneg\phi(f,s_1,...,s_n)\rbc$, and this is a 
contradiction.

\end{proof}

\end{document}
